\section{Análisis Descriptivo (EDA)}
\label{sec:descriptivo}

\subsection*{Estadísticos generales de la variable objetivo}

\begin{minted}{python}
power_cols = [c for c in df.columns if 'ac_power' in c]
df['total_ac_power'] = df[power_cols].sum(axis=1, min_count=1)
df['total_ac_power'].describe()
\end{minted}

\begin{table}[H]
\centering
\begin{tabular}{lr}
\toprule
\textbf{Estadístico} & \textbf{total\_ac\_power (W)} \\
\midrule
count & 632\,952 \\
media & 156.32 \\
desv. estándar & 230.96 \\
mínimo & 0.00 \\
percentil 25 & 0.00 \\
mediana & 0.00 \\
percentil 75 & 291.60 \\
máximo & 721.61 \\
\bottomrule
\end{tabular}
\caption{Estadísticos descriptivos de la potencia AC total de la planta}
\end{table}

Como es característico de la generación solar, la mediana de \texttt{total\_ac\_power} es 0: más de la mitad de los registros corresponden a horas sin generación (nocturnas o con inversores fuera de línea), mientras que el 25\% superior de los registros supera los 291.6 W, reflejando picos de generación cercanos al mediodía. El porcentaje de lecturas en cero por inversor confirma este patrón y varía entre 52.5\% y 59.2\% según el inversor, ligeramente por encima del $\sim$50\% esperado solo por el ciclo día/noche —diferencia que se explica por los periodos adicionales de inactividad de cada equipo (Sección \ref{sec:no_supervisado}).

\subsection*{Distribuciones univariadas}

\begin{figure}[H]
\centering
\includegraphics[width=\textwidth]{fig_distribuciones.png}
\caption{Distribuciones de la potencia AC total de la planta, la corriente DC y el voltaje AC de un inversor de referencia.}
\end{figure}

La potencia AC total muestra una masa importante en cero (rango nocturno) y una cola derecha correspondiente a las horas de mayor irradiancia. El voltaje AC, en contraste, se concentra en un rango angosto: es una variable regulada por el inversor y varía poco entre condiciones de operación, a diferencia de la corriente, que sí responde directamente a la irradiancia disponible.

\subsection*{Comportamiento por dispositivo}

\begin{figure}[H]
\centering
\includegraphics[width=\textwidth]{fig_boxplot_inversores.png}
\caption{Distribución de la potencia AC por inversor (solo horas con generación $>0$).}
\end{figure}

La mediana de potencia por inversor (calculada solo sobre horas con generación) oscila entre 8.9 W (inversor 07) y 13.9 W (inversor 20). Los inversores \textbf{04} y \textbf{07} destacan tanto por su mediana más baja como por presentar valores máximos anómalos ($\sim$104--105 W, frente a $\sim$30 W del resto de la flota) — un indicio de comportamiento atípico que se confirma cuantitativamente en el análisis de clustering (Sección \ref{sec:no_supervisado}).

\subsection*{Patrones temporales}

\begin{figure}[H]
\centering
\includegraphics[width=\textwidth]{fig_perfiles_temporales.png}
\caption{Perfil diurno promedio (izquierda) y perfil mensual promedio (derecha) de la potencia AC total generada.}
\end{figure}

El perfil diurno confirma la forma de campana esperada en generación solar (nula de noche, máxima cerca del mediodía). El perfil mensual refleja la estacionalidad anual asociada a la posición del sol y a los patrones de nubosidad de cada época del año.

\begin{figure}[H]
\centering
\includegraphics[width=\textwidth]{fig_serie_diaria.png}
\caption{Potencia AC total de la planta — promedio diario, 2017--2023.}
\end{figure}

La serie diaria completa permite observar tanto la estacionalidad anual repetida a lo largo de los $\sim$6 años de historia como posibles cambios de nivel (p. ej. mantenimiento, expansión o degradación de la planta) que serán relevantes para el forecasting (Sección \ref{sec:forecasting}).
